\subsection{Regular languages and expressions}

\begin{definition}
	Let $\Sigma$ be an alphabet. The class of \emph{regular languages} over
	$\Sigma$ is the inclusion-wise minimal subset of $\Sigma^*$ such that
	\begin{itemize}
	\item $\emptyset$ is a regular language over $\Sigma$,
	\item if $\sigma \in \Sigma_\epsilon$, then $\{ \sigma \}$ is a regular
		language over $\Sigma$,
	\item if $\mathcal{L}_1$ and $\mathcal{L}_2$ are regular languages over
		$\Sigma$, then
		\begin{itemize}
		\item $\mathcal{L}_1 \hat{\cdot} \mathcal{L}_2$ is a regular
			language over $\Sigma$,
		\item $\mathcal{L}_1 \cup \mathcal{L}_2$ is a regular language over
			$\Sigma$,
		\item $\mathcal{L}_1^*$ is a regular language over $\Sigma$.
		\end{itemize}
	\end{itemize}
\end{definition}

\begin{definition}
	Let $\Sigma$ be an alphabet such that $\{ \emptyset, (, ), ., +, * \}
	\not\subseteq \Sigma$. The class of \emph{regular expressions} over $\Sigma$
	is the inclusion-wise minimal subset of $(\Sigma \cup \{ \emptyset, (, ),
	., +, * \})^*$ such that
	\begin{itemize}
	\item $\emptyset$ is a regular expression over $\Sigma$,
	\item if $\sigma \in \Sigma_\epsilon$, then $\sigma$ is a regular expression
		over $\Sigma$,
	\item if $\mathcal{R}_1$ and $\mathcal{R}_2$ are regular expressions over
		$\Sigma$, then
		\begin{itemize}
		\item $(\mathcal{R}_1 . \mathcal{R}_2)$ is a regular expression over
			$\Sigma$,
		\item $(\mathcal{R}_1 + \mathcal{R}_2)$ is a regular expression over
			$\Sigma$,
		\item $(\mathcal{R}_1*)$ is a regular expression over $\Sigma$.
		\end{itemize}
	\end{itemize}
\end{definition}

\begin{definition}
	The \emph{language denoted by a regular expression} $\mathcal{R}$ over an
	alphabet $\Sigma$ is the set $\mathcal{L}_\mathcal{R}$ defined as
	\begin{itemize}
	\item $\emptyset$ if $\mathcal{R} = \emptyset$,
	\item $\{ \sigma \}$ if $\mathcal{R} = \sigma$ for $\sigma \in
		\Sigma_\epsilon$,
	\item $\mathcal{L}_{\mathcal{R}_1} \hat{\cdot} \mathcal{L}_{\mathcal{R}_2}$
		if $\mathcal{R} = (\mathcal{R}_1 . \mathcal{R}_2)$,
	\item $\mathcal{L}_{\mathcal{R}_1} \cup \mathcal{L}_{\mathcal{R}_2}$
		if $\mathcal{R} = (\mathcal{R}_1 + \mathcal{R}_2)$,
	\item $\mathcal{L}_1^*$ if $\mathcal{R} = (\mathcal{R}_1*)$.
	\end{itemize}
\end{definition}

\begin{theorem}[Thompson, see \cite{T68}]\label{thm:thompson}
	Let $\mathcal{R}$ be a regular expression over an alphabet $\Sigma$. There
	exists a finite-state automaton $\mathcal{A}$ such that
	$\mathcal{L}_\mathcal{A} = \mathcal{L}_\mathcal{R}$.
\end{theorem}

The proof of \Cref{thm:thompson} is essentially a recursive method for
constructing a finite-state automaton $\mathcal{A}$ from a regular expression
$\mathcal{R}$. It effectively creates 5 distinct types of finite-state automata
depending on the type of regular expression.
\begin{itemize}
\item If $\mathcal{R} = \emptyset$, then $\mathcal{A} = \langle \Sigma, \{ s, f \},
	\{ s \}, \{ f \}, \emptyset \rangle$ is the Type 1 finite-state automaton
	constructed from $\mathcal{R}$,
\item if $\mathcal{R} = \sigma$ for $\sigma \in \Sigma_\epsilon$, then $\mathcal{A}
	= \langle \Sigma, \{ s, f \}, \{ s \}, \{ f \}, \{ \langle s, \sigma, f
	\rangle \} \rangle$ is the Type 2 finite-state automaton constructed from
	$\mathcal{R}$,
\item if $\mathcal{A}' = \langle \Sigma, Q', S', F', \Delta' \rangle$ and
	$\mathcal{A}'' = \langle \Sigma, Q'', S'', F'', \Delta'' \rangle$ are the
	finite-state automata constructed respectively from $\mathcal{R}'$ and
	$\mathcal{R}''$, then
	\begin{itemize}
	\item $\mathcal{A} = \langle \Sigma, Q' \cup Q'', S', F'', \Delta \rangle$
		is the Type 3 finite-state automaton constructed from $(\mathcal{R}'
		. \mathcal{R}'')$ where $\Delta = \Delta' \cup \Delta'' \cup \{
		\langle f', \epsilon, s'' \rangle \mid \exists f' \in F' \; \exists
		s'' \in S'' \}$,
	\item $\mathcal{A} = \langle \Sigma, Q' \cup Q'' \cup \{ s, f \}, \{ s \},
		\{ f \}, \Delta \rangle$ is the Type 4 finite-state automaton
		constructed from $(\mathcal{R}'+\mathcal{R}'')$ where $\Delta =
		\Delta' \cup \Delta'' \cup \{ \langle s, \epsilon, s' \rangle,
		\langle s, \epsilon, s'' \rangle, \langle f', \epsilon, f \rangle,
		\langle f'', \epsilon, f \rangle \mid \exists s' \in S' \; \exists
		s'' \in S'' \; \exists f' \in F' \; \exists f'' \in F'' \}$,
	\item $\mathcal{A} = \langle \Sigma, Q' \cup \{ s, f \}, \{ s \}, \{ f \},
		\Delta \rangle$ is the Type 5 finite-state automaton constructed
		from $(\mathcal{R}'*)$ where $\Delta = \Delta' \cup \{ \langle s,
		\epsilon, s' \rangle, \langle s, \epsilon, f \rangle$, $\langle f',
		\epsilon, f \rangle, \langle f', \epsilon, s' \rangle \mid \exists
		s' \in S' \; \exists f' \in F' \}$.
	\end{itemize}
\end{itemize}

\begin{figure}[!h]
	\centering
	\begin{tikzpicture}[on grid, auto, scale=0.35]
  		\node[state, initial, fill=white, scale=0.375]
			(s1) at (-10, 0) {};
  		\node[state, fill=white, scale=0.375]
			(s2) at (-6, 0) {};
		\node[state, fill=white, scale=0.375]
			(s3) at (-2, 2) {};
		\node[state, fill=white, scale=0.375]
			(s4) at (-2, -2) {};
		\node[state, fill=white, scale=0.375]
			(s5) at (2, 2) {};
		\node[state, fill=white, scale=0.375]
			(s6) at (2, -2) {};
		\node[state, fill=white, scale=0.375]
			(s7) at (6, 0) {};
  		\node[state, accepting, fill=white, scale=0.375]
			(s8) at (10, 0) {};

		\draw
			(s1) edge node{$\epsilon$} (s2)
			(s2) edge node{$\epsilon$} (s3)
			(s2) edge node{$\epsilon$} (s4)
			(s3) edge node{$0$} (s5)
			(s4) edge node{$1$} (s6)
			(s5) edge node{$\epsilon$} (s7)
			(s6) edge node{$\epsilon$} (s7)
			(s7) edge node{$\epsilon$} (s8)
			(s1) edge[bend right=40] node{$\epsilon$} (s8)
			(s7) edge[bend left=-80] node[above]{$\epsilon$} (s2);
	\end{tikzpicture}
	\caption{Thompson finite-state automaton derived from the regular expression
		$(0+1)*$  by applying the construction described in
		\Cref{thm:thompson}.}
	\label{fig:f_thompson_0}
\end{figure}

\begin{proposition}
	Let $\mathcal{A} = \langle \Sigma, Q, S, F, \Delta \rangle$ be a
	finite-state automaton derived from a regular expression by applying the
	construction described in \Cref{thm:thompson}. Then $|S| = |F| = 1$.
\end{proposition}

\begin{proof}
	We prove the statement by induction on the construction of the finite-state
	automaton.
	\begin{case}[$\mathcal{A}$ is Type 1, 2, 4 or 5]
	$|S| = |\{ s \}| = 1$ and $|F| = |\{ f \}| = 1$ by construction.
	\end{case}
	\begin{case}[$\mathcal{A}$ is Type 3]
	$|S| = |S'| = 1$ and $|F| = |F''| = 1$ by the inductive hypothesis for
	$\mathcal{A'}$ and $\mathcal{A''}$.
	\end{case}
\end{proof}
